\section{Program and Firmware Flow}
\label{sec:ch3-program-flow}

\subsection{Initialization and Sampling}
\label{sec:ch3-initialization-sampling}

The firmware follows a repeated acquisition, computation, inference, and
protection cycle. On startup, the ESP32-S3 initializes the sensors, relay
driver, OLED display, buzzer, Wi-Fi, Firebase interface, and local protection
state. In the main loop, it samples the voltage, thermistor, and current
waveform, applies calibration and filtering, computes thermal and waveform
features, updates the context classifier, evaluates the arc classifier, and
passes the resulting variables to the protection manager.

Figure~\ref{fig:ch3-program-flow} shows the operating sequence implemented by
the firmware.
\begin{figure}[H]
  \centering
  \includegraphics[width=0.96\textwidth]{ch3/figure_3.3.png}
  \caption{Program flowchart of the TinyML Smart Plug firmware.}
  \label{fig:ch3-program-flow}
\end{figure}

\subsection{Protection-State Evaluation}
\label{sec:ch3-protection-state-evaluation}

Figure~\ref{fig:ch3-program-flow} shows that a relay trip can be caused by a
sustained overload, an estimated socket over-temperature, or an arc-fault
classification. The overload path issues a warning above 10.0A, but relay
disconnection is based on persistence above 14.5A. This persistence is
implemented as a score that increases every 10s while the
condition remains active and trips when the score reaches 6, equivalent to
approximately 60s of sustained overload. This makes the overload behavior
closer to inverse-time protection logic than to an instantaneous cutoff [22].

The firmware also latches fault states. After relay disconnection, the event is
recorded locally and uploaded when the network is available. The plug remains
in the fault state until the user clears the fault through the device or
application workflow. This prevents a hazardous load from being automatically
re-energized after a transient communication or reset event.

\subsection{Timing Constants}
\label{sec:ch3-timing-constants}

The sampling and timing constants used by the firmware are listed in
Table~\ref{tab:ch3-firmware-constants}. These values were taken from the
firmware configuration and the current-sensing implementation.
\begin{table}[H]
  \centering
  \caption{Selected firmware timing and sampling constants}
  \label{tab:ch3-firmware-constants}
  \begin{tabular}{p{0.38\textwidth}p{0.22\textwidth}p{0.30\textwidth}}
    \toprule
    \textbf{Constant} & \textbf{Value} & \textbf{Role} \\
    \midrule
    Target current sampling rate & 27,000Hz & Current waveform acquisition for feature extraction \\
    Arc runtime frame length & 512 samples & FFT and time-domain feature window \\
    Arc runtime hop length & 256 samples & Overlap between consecutive feature frames \\
    Feature emission cadence & Every 3 hops & Approximate 35.16Hz feature output rate \\
    MCP3204 SPI clock & 1.2MHz & 12-bit current ADC transfer clock \\
    Voltage sampling interval & 250 $\mu$s & Firmware voltage-window acquisition interval \\
    \bottomrule
  \end{tabular}
\end{table}

Table~\ref{tab:ch3-firmware-constants} explains why the current path is
handled separately from slow telemetry. Arc-fault detection requires waveform
frames with millisecond-scale timing, while Firebase updates and PWA displays
operate at a slower human-monitoring cadence.
